\section{Background and Motivation}
\label{sec:background}

\subsection{Agent tool workflows}

An agent is a loop around a language model. In each round the model receives
the task and everything observed so far, and emits either a \emph{tool call}
or a statement that it is done. A tool call reads or writes shared state: a
deployment, a calendar, a customer record. The result of the call is appended
to the model's context, and the loop continues.

Two properties of this loop matter for concurrency control. First, the
interval between two tool calls is a model inference. It takes seconds, it is
private to the agent, and it touches no shared state. We call it the
\emph{reasoning gap}. Second, the sequence of calls is not known in advance:
which tool the agent calls next depends on what the last call returned.

When several agents work on the same system, their tool calls interleave. The
result is the familiar catalogue of anomalies~\cite{bernstein1987}. Two agents
that both read a value, reason, and write a derived value lose an update. An
agent that scans a collection and then acts on what it saw misses an object
another agent inserted in between, a phantom~\cite{eswaran1976}. An agent that
appends a record derived from a value another agent has since overwritten
leaves a log that no serial execution could have produced.

\subsection{Existing answers}

\paragraph{Whole-workflow isolation.} Treating each agent's entire workflow
as one transaction prevents every anomaly above. Under strict two-phase
locking (2PL) the agent holds each lock from first use until it finishes;
under optimistic concurrency control (OCC)~\cite{kung1981occ} it runs
unlocked and validates at the end. Because the target systems that agents
operate on---a Kubernetes API, a CRM---offer neither private workspaces nor
rollback, writes land in place, and an abort is implemented as a
saga~\cite{garciamolina1987sagas}: each tool registers an inverse, and a
victim's writes are undone by running the inverses. The cost of both policies
is tied to the reasoning gap. 2PL holds locks across it; OCC discards a whole
workflow's worth of inference when validation fails.

\paragraph{Ordered repair.} CoAgent~\cite{lyu2026coagent} takes the opposite
position. Its protocol, \mtpo, assigns every agent a rank $\sigma$ at launch
and aims for the outcome of running the agents serially in $\sigma$ order.
Reads return the value the reader would have seen in that order; writes apply
speculatively in place; and when a lower-ranked agent overwrites something a
higher-ranked agent has read, the higher-ranked agent receives a
\emph{notification} carrying the new value. Notifications flow in one
direction only, from lower $\sigma$ to higher. Nobody blocks and nothing is
undone, which makes \mtpo fast. Its correctness, however, rests on two
assumptions about the model: that an agent completes its task correctly in
isolation (A1), and that a notified agent correctly judges whether its plan is
affected and repairs exactly the affected operations (A3). CoAgent assumes
A3 holds 95\% of the time. We return to both assumptions in
\S\ref{sec:eval:agent}.

\subsection{Transaction chopping}

Transaction chopping~\cite{shasha1995chopping} asks a different question:
which transactions can be cut into smaller ones without anyone being able to
tell? A \emph{chopping} splits each transaction into consecutive
\emph{pieces}. Each piece executes as its own serializable transaction and
commits before the next begins, so locks are held only within a piece.

Whether a chopping is safe is decided by an undirected graph over the pieces,
the \emph{SC graph}. Two pieces of the same transaction are joined by an
\emph{S~edge} (sibling). Two pieces of different transactions that access a
common item, at least one of them writing it, are joined by a \emph{C~edge}
(conflict). The chopping is correct---every execution is equivalent to a
serial execution of the \emph{original} transactions---if the graph has no
simple cycle containing both an S~edge and a C~edge.

\begin{figure}[t]
\centering
\begin{tikzpicture}[
  piece/.style={draw=ccmuted, rounded corners=2pt, minimum height=5mm,
                minimum width=11mm, font=\scriptsize, inner sep=1.5pt},
  S/.style={ccblue, line width=0.9pt},
  C/.style={ccorange, line width=0.9pt, dashed},
  lab/.style={font=\scriptsize}]
% ---- (a) choppable ------------------------------------------------------
\node[lab] at (1.05,1.75) {(a) certifiable};
\node[lab, anchor=east] at (-0.1,1.1) {$A$};
\node[lab, anchor=east] at (-0.1,0.0) {$B$};
\node[piece] (a1) at (0.5,1.1) {$r(\mathit{pool})$};
\node[piece] (a2) at (1.9,1.1) {$w(c_A)$};
\node[piece] (b1) at (0.5,0.0) {$w(\mathit{pool})$};
\node[piece] (b2) at (1.9,0.0) {$w(c_B)$};
\draw[S] (a1) -- (a2);
\draw[S] (b1) -- (b2);
\draw[C] (a1) -- (b1);
% ---- (b) not choppable --------------------------------------------------
\begin{scope}[xshift=4.0cm]
\node[lab] at (1.05,1.75) {(b) mixed cycle};
\node[lab, anchor=east] at (-0.1,1.1) {$A$};
\node[lab, anchor=east] at (-0.1,0.0) {$B$};
\node[piece] (c1) at (0.5,1.1) {$r(x)$};
\node[piece] (c2) at (1.9,1.1) {$w(x)$};
\node[piece] (d1) at (0.5,0.0) {$r(x)$};
\node[piece] (d2) at (1.9,0.0) {$w(x)$};
\draw[S] (c1) -- (c2);
\draw[S] (d1) -- (d2);
\draw[C] (c1) -- (d2);
\draw[C] (c2) -- (d1);
\draw[C] (c2) -- (d2);
\end{scope}
% ---- legend ---------------------------------------------------------------
\draw[S] (0.6,-0.75) -- (1.2,-0.75) node[right, lab, black] {S edge (same workflow)};
\draw[C] (3.9,-0.75) -- (4.5,-0.75) node[right, lab, black] {C edge (conflict)};
\end{tikzpicture}
\caption{SC graphs for two agent workflows cut at the reasoning gap. (a)~$A$
reads a shared pool and writes a private claim while $B$ updates the pool: the
only conflict joins the two first pieces, no cycle mixes edge types, and the
chopping is safe. (b)~Both agents read and then write the same object: the
cycle $r_A(x)$--$w_A(x)$--$w_B(x)$--$r_A(x)$ mixes an S~edge with C~edges, so
the cut would permit a lost update and each workflow must stay one piece.}
\label{fig:sc}
\end{figure}

Figure~\ref{fig:sc} shows the two cases that matter for agents. When the
shared dependency is short and the work that follows is private, cutting at
the reasoning gap is safe, and the lock on the shared item is released before
the model starts thinking. When a read-modify-write straddles the gap, the
graph has a mixed cycle and the cut must not be made. The repair is always
available before execution: merge pieces. A workflow that is a single piece
has no S~edges, so it cannot be on a mixed cycle.

\subsection{Why agents break the classical setting}

The chopping theorem quantifies over a set of transaction programs that is
fixed before any piece runs. An agent's plan is not fixed. Suppose a workflow
is certified as ``read $x$; then write a private result'', the read commits,
and the agent then decides, on the basis of what it read, that it must also
write a shared item $y$. The certified plan no longer describes the execution.

Two responses are unsound or unavailable. Certifying only the new suffix is
unsound, because the conflicts that close a cycle may run through the pieces
that already committed. Merging the new write into the piece that performed
the read is unavailable, because that piece has committed and released its
locks; other workflows may already have observed its effects or interleaved
with it. What remains is to regroup the \emph{uncommitted} part of the plan,
and the question this paper answers is when that is enough.
